% Generated draft: requires mathematical and source review.
\documentclass[11pt]{article}
\usepackage[margin=25mm]{geometry}
\usepackage{fontspec}
\setmainfont{Latin Modern Roman}
\usepackage{amsmath,amssymb,mathtools,unicode-math}
\setmathfont{Latin Modern Math}
\usepackage{enumitem,needspace,xurl,hyperref}
\setlength{\parindent}{0pt}
\setlength{\parskip}{5pt}
\setlength{\emergencystretch}{3em}
\newcommand{\N}{\mathbb{N}}
\newcommand{\Z}{\mathbb{Z}}
\newcommand{\Q}{\mathbb{Q}}
\newcommand{\R}{\mathbb{R}}
\newcommand{\C}{\mathbb{C}}
\newcommand{\F}{\mathbb{F}}
\newcommand{\A}{\mathbb{A}}
\newcommand{\PP}{\mathbb{P}}
\newcommand{\E}{\mathbb{E}}
\newcommand{\eps}{\varepsilon}
\newcommand{\dd}{\,\mathrm{d}}
\begin{document}
\section*{UnsolvedMath ID: 30001297}
\textbf{Source problem code:} catalogue entry
\subsection{Definitions and mathematical statement}
Let \(k\) be a number field, let \(\bar k\) be an algebraic closure, and let \(\Omega_k\) denote the set of all places of \(k\). A K3 surface over \(k\) is a smooth, projective, geometrically integral surface \(X\) such that
\[
\omega_{X_{\bar k}}\cong\mathcal O_{X_{\bar k}},\qquad H^1(X_{\bar k},\mathcal O_{X_{\bar k}})=0.
\]
For \(v\in\Omega_k\), write \(k_v\) for the completion of \(k\). Since \(X\) is proper, its adelic points may be identified with
\[
X(\mathbb A_k)=\prod_{v\in\Omega_k}X(k_v),
\]
with the usual adelic topology. Write
\[
\operatorname{Br}(X)=H^2_{\mathrm{\acute et}}(X,\mathbb G_m)_{\mathrm{tors}}
\]
for its cohomological Brauer group. For \(\alpha\in\operatorname{Br}(X)\) and \(P_v\in X(k_v)\), let \(\alpha(P_v)\in\operatorname{Br}(k_v)\) be its evaluation, and let \(\operatorname{inv}_v:\operatorname{Br}(k_v)\to\mathbb Q/\mathbb Z\) be the local invariant map. The Brauer--Manin set is
\[
X(\mathbb A_k)^{\operatorname{Br}}=
\left\{(P_v)_v\in X(\mathbb A_k):
\sum_{v\in\Omega_k}\operatorname{inv}_v(\alpha(P_v))=0
\text{ for every }\alpha\in\operatorname{Br}(X)\right\}.
\]
Each such sum has only finitely many nonzero terms. Global reciprocity places the diagonal image of \(X(k)\) in this set.

\textbf{Conjecture (as supplied).} For every number field \(k\) and every K3 surface \(X/k\), the diagonal image of \(X(k)\) is dense in \(X(\mathbb A_k)^{\operatorname{Br}}\). Explicitly, for every \((P_v)_v\in X(\mathbb A_k)^{\operatorname{Br}}\), every finite subset \(S\subset\Omega_k\), and every choice of open neighborhoods \(U_v\subset X(k_v)\) of \(P_v\) for \(v\in S\), there exists \(P\in X(k)\) whose image lies in \(U_v\) for every \(v\in S\).

In particular, nonemptiness of the Brauer--Manin set would imply \(X(k)\ne\varnothing\), and every finite local approximation arising from a Brauer--Manin adelic point would be realizable by a rational point. These are the supplied meanings of the Brauer--Manin obstruction being the only obstruction to the Hasse principle and weak approximation.

The formulation above uses the full adelic topology, including the usual topology at archimedean places. The supplied record does not specify whether any modification at archimedean places was intended; that source convention requires review.
\subsection{Short English statement}
The supplied conjecture asks whether, for every K3 surface over every number field, rational points can approximate every adelic point orthogonal to the entire Brauer group, at any finitely many places simultaneously. Its supplied open status is dated curation, not an independently verified conclusion.
\subsection{Sources}
\begin{itemize}
\item Supplied local UnsolvedMath registry record, permanent ID \(30001297\), title ``Brauer--Manin Conjecture for K3 Surfaces,'' source problem code \texttt{catalogue entry}, in the Oberwolfach Reports Open Problems collection. No individual proposer is supplied; the recorded proposed year is 2009.
\item Supplied report citation: \emph{Explicit Methods in Number Theory (2010)}, Oberwolfach Reports, DOI 10.4171/owr/2009/33, pp. 1843--1920. Source URL and literature-source link labeled ``OWR report'': \url{https://doi.org/10.4171/owr/2009/33}. The supplied background describes these reports as documenting mathematical workshops and their open-problem discussions.
\item Supplied literature-source link labeled ``Skorobogatov on Brauer groups and rational points'': \url{https://doi.org/10.1090/surv/198}. No more detailed bibliographic identification or pinpoint reference is supplied.
\item Literature assessment checked on \textbf{2026-08-21}, according to the supplied record: status ``open''; density in the Brauer--Manin set for K3 surfaces over number fields is described as conjectural in general, with known results treating special families and obstructions. This is supplied dated curation derived from the OWR material, not independently verified literature assessment. The supplied background expressly requests independent verification of time-sensitive status and cited claims.
\item The registry labels the statement \texttt{corrected\_verified} and describes its editing as meaning-preserving semantic, typographic, and LaTeX cleanup without unsupported source-dependent additions. These are supplied editorial assertions, not independent source verification or mathematical review.
\end{itemize}
\end{document}
